\begin{frame}
    \frametitle{Math}
    \framesubtitle{Syntax}
    \begin{center}
    \begin{tabular}{ccl}
        \textbf{Environment}
            & \textbf{Symbol(s)}
            & \textbf{Description}\\
    \toprule[1.2pt]
        \lstinline|math|
            & \lstinline|$...$| 
            & Inline math delimiters\\
    \midrule\pause
        \lstinline|displaymath|
            & \lstinline|\\[...\\]|
            & Display math delimiters\\
    \end{tabular}
    \end{center}
\end{frame}

\begin{frame}
    \frametitle{Math}
    \framesubtitle{Operator Symbols}
    \begin{center}
    \begin{tabular}{l c c >{\raggedright\arraybackslash}m{5cm}}
        \textbf{Operation}
        & \textbf{Command}
        & \textbf{Example}
        & \textbf{Notes}\\
    \toprule[1.2pt]\pause
        Addition
        & \texttt{+}
        & $a + b$
        & \\
    \midrule\pause
        Subtraction
        & \texttt{-}
        & $c - d$
        & \\
    \midrule\pause
        Multiplication
        & \lstinline|\\cdot|
        & $e \cdot f$
        & Do \emph{not} use \texttt{*} or \lstinline|\\times|!\\
    \midrule\pause
        Division
        & \lstinline|\\frac\{\}\{\}|
        & $\dfrac{g}{h}$ % \dfrac{}{} used here for better visualization
        & For inline math a slash \texttt{/} is also acceptable.\\
    \midrule\pause
        Exponential
        & \lstinline|^\{\}|
        & $i^{j}$
        & Sometimes also referred to as \emph{superscript}.\\
    \midrule\pause
        Root
        & \lstinline|\\sqrt[]\{\}|
        & $\sqrt{k + l}$
        & The optional parameter is for non-square roots e.g. $\sqrt[3]{m}$.\\
    \end{tabular}
    \end{center}
\end{frame}

\begin{frame}
    \frametitle{Math}
    \framesubtitle{Text In Math and Indices}
    \only<1-4|handout:1>{%
    \begin{center}
    \begin{tabular}{c p{18em}}
        \textbf{Command}
        & \textbf{Definition}\\
    \toprule[1.2pt]\pause
        \lstinline|\\text\{\}|
        & Inserts regular text\\
    \midrule\pause
        \lstinline|\\mathrm\{\}|
        & Sets font to math-roman (\emph{not} equivalent to \lstinline|\\text\{\}|)\\
    \midrule\pause
        \lstinline|_\{\}|
        & Indices of a variable. Sometimes also referred to as \emph{subscript}.\\
    \end{tabular}
    \end{center}}
    \only<5-|handout:2>{%
    \begin{columns}
        \begin{column}{.54\textwidth}
            \tcbexinputlst{listings/math-text.tex}
        \end{column}
        \begin{column}{.44\textwidth}
            \[
  V_{\mathrm{out}}
  = \begin{cases}
    0 &\text{if} \quad t < 0\\
    1 &\text{else}
  \end{cases}
\]
        \end{column}
    \end{columns}
    }
\end{frame}

\begin{frame}
    \frametitle{Math}
    \framesubtitle{Functions}
    \begin{center}
    \begin{tabular}{l c c}
        \textbf{Function}
        & \textbf{Command}
        & \textbf{Example}\\
    \toprule[1.2pt]\pause
        Sine
        & \lstinline|\\sin|
        & $\sin\alpha$\\
    \midrule\pause
        Cosine
        & \lstinline|\\cos|
        & $\cos\beta$\\
    \midrule\pause
        Tangens
        & \lstinline|\\tan|
        & $\tan\gamma$\\
    \end{tabular}
    \end{center}
\end{frame}

% TODO: add link/reference to https://oeis.org/wiki/List_of_LaTeX_mathematical_symbols

\begin{frame}
    \frametitle{Math}
    \begin{columns}
        \begin{column}{.54\textwidth}
            \tcbexinputlst{listings/math-cosine-law.tex}
        \end{column}
        \begin{column}{.44\textwidth}
            \input{listings/math-cosine-law.tex}
        \end{column}
    \end{columns}
\end{frame}

\begin{frame}
    \frametitle{Math}
    \begin{columns}
        \begin{column}{.54\textwidth}
            \tcbexinputlst{listings/math-sine-law.tex}
        \end{column}
        \begin{column}{.44\textwidth}
            The law of sines relates triangle side lengths to its angles with
\[
  \frac{a}{\sin\alpha}
  =
  \frac{b}{\sin\beta}
  =
  \frac{c}{\sin\gamma}.
\]
        \end{column}
    \end{columns}
\end{frame}

\begin{frame}
    \frametitle{Math}
    \framesubtitle{More functions}
    \begin{center}
    \begin{tabular}{l c c}
        \textbf{Function}
        & \textbf{Command}
        & \textbf{Example}\\
    \toprule[1.2pt]\pause
        Limit
        & \lstinline|\\lim_\{\}|
        & $\lim_{n \to \infty} \frac{1}{n}$\\
    \midrule\pause
        Sum
        & \lstinline|\\sum_\{\}^\{\}|
        & $\sum_{i = 0}^{n} a_n$\\
    \midrule\pause
        Product
        & \lstinline|\\prod_\{\}^\{\}|
        & $\prod_{j = 0}^{m}b_m$\\
    \end{tabular}
    \end{center}
\end{frame}

\begin{frame}
    \frametitle{Math}
    \framesubtitle{Integrals and Derivatives}
    \only<1-3>{%
    \begin{center}
    \begin{tabular}{l c c}
        \textbf{Operation}
        & \textbf{Command}
        & \textbf{Example}\\
    \toprule[1.2pt]\pause
        Derivative
        & \lstinline|\\frac\{d\\,\}\{d\\, x\}|
        & $\frac{d\,f(x)}{d\, x}$\\
    \midrule\pause
        Integral
        & \lstinline|\\int_\{\}^\{\} d\\,|
        & $\int_{a}^{b} x d\,x$\\
    \end{tabular}
    \end{center}
    }
    \begin{alertblock}<4->{\blockstrut Remark}
        Some people prefer typesetting it as
        \[
            \frac{\operatorname{d} }{\operatorname{d} x}
            \qquad
            \text{and}
            \qquad
            \int x \operatorname{d} x,
        \]
        using \lstinline|\\operatorname\{d\}| instead of simply \lstinline|d|.
    \end{alertblock}
\end{frame}